\section{Conclusion}\label{sec:conclusion}
Neural Bellman Operators connect policy evaluation, feasible improvement, and independent economic error accounting. The capital application now distinguishes an architectural performance envelope from a policy-specific post-training certificate. A paired payoff identity, explicit diffusion-transfer estimate, and finite-family sampling bound identify the value of a deployed neural correction relative to an analytical schedule. The economic comparison remains the original continuous problem and its full adapted policy class.

The rollout-costate implementation evaluates the current policy before improving its Hamiltonian. Direct-policy, affine, and tuned critic-greedy comparisons expose the cost and accuracy of that implementation without presuming its superiority. A positive verified gain is evidence about a particular deployed policy; it is neither an optimizer convergence theorem nor a claim about every initialization. The reported improvement is a quantitative result in the stated stylized economy, not an empirical calibration.

Recursive utility, endogenous preferences, sophisticated temporal selves, dynamic games, viscosity selection, and stochastic traces remain within the paper's economic program. Their full equations and evidence are retained in the supplement, with their own domains and verification hypotheses. Extending the new diffusion-transfer calculation to other observation rules, state-dependent diffusion, recursive utility, or stopping boundaries requires those additional error terms to be verified. The present theorem does not silently supply them.

